\exercisesection{}

\begin{enumerate}
\def\labelenumi{\arabic{enumi})}
\tightlist
\item
  设 \(\mathbf{E}\) 为一个集合，\(\Phi\) 为 \(\mathbf{E}\) 的子集所成的一个集环，\(m\) 为定义在 \(\Phi\) 上的一个可加集函数。用 \(\mathcal{E}(\Phi)\) 表示实 \(\Phi\) -阶梯函数在 \(\mathbf{E}\) 上所成的向量空间，并用 \(J\) 表示 \(\mathcal{E}(\Phi)\) 上满足 \(J(\varphi_{\mathbf{A}}) = m(\mathbf{A})\)（对所有 \(A \in \Phi\)）的线性形式。
\end{enumerate}

\begin{enumerate}
\def\labelenumi{\alph{enumi})}
\tightlist
\item
  对每个 \(\mathbf{A} \in \Phi\)，置
\end{enumerate}

\[
m _ {+} (\mathrm{A}) = \sup _ {\mathrm{B} \in \Phi , \mathrm{B} \subset \mathrm{A}} m (\mathrm{B}) \quad \text { and } \quad m _ {-} (\mathrm{A}) = \sup _ {\mathrm{B} \in \Phi , \mathrm{B} \subset \mathrm{A}} (- m (\mathrm{B})).
\]

证明 \(|m|(\mathbf{A}) = m_{+}(\mathbf{A}) + m_{-}(\mathbf{A})\) 是形如 \(\sum_{i=1}^{p}|m(\mathbf{A}_i)|\) 的数所成集合的上确界（对所有有限划分 \((\mathbf{A}_i)_{1\leqslant i\leqslant p}\)（把 A 划分成属于 \(\Phi\) 的集合））。

\begin{enumerate}
\def\labelenumi{\alph{enumi})}
\setcounter{enumi}{1}
\item
  称 \(m\) 是相对有界的，如果形如 \(m(\mathbf{B})\)（其中 \(\mathbf{B} \in \Phi\)，\(\mathbf{B} \subset \mathbf{A}\)）的数所成的集合对任何集合 \(\mathbf{A} \in \Phi\) 都是有界的。证明下列条件等价：\(\alpha)\) \(m\) 相对有界；\(\beta)\) 线性形式 \(J\) 在 Riesz 空间 \(\mathcal{E}(\Phi)\) 上相对有界；\(\gamma)\) \(m\) 是 \(\Phi\) 上两个正可加集函数之差；\(\delta)\) \(|m|(\mathbf{A})\) 对所有 \(\mathbf{A} \in \Phi\) 有限（按 \(\alpha) \Rightarrow \beta) \Rightarrow \gamma) \Rightarrow \delta) \Rightarrow \alpha)\) 循环论证）。
\item
  假设 \(m\) 相对有界。证明集函数 \(|m|\)、\(m_{+}\) 与 \(m_{-}\) 都是可加的，且 \(m = m_{+} - m_{-}\)（利用 Ch. II, §2, No.~2 的 Th. 1）。此外，若 \(m'\) 与 \(m''\) 是 \(\Phi\) 上满足 \(m = m' - m''\) 的可加正集函数，则 \(m' \geqslant m_{+}\) 且 \(m'' \geqslant m_{-}\)。
\end{enumerate}

\(d)\) 称 \(m\) 是可数可加的，如果 \(m(A) = \sum_{n\geqslant 1}m(A_n)\) 对每个可数

划分 \((\mathbf{A}_n)_{n\geqslant 1}\)（把集合 \(\mathbf{A}\in \Phi\) 划分成属于 \(\Phi\) 的集合）都成立。证明这一条件意味着：\(\lim_{n\to \infty}m(\mathbf{A}_n) = 0\) 对每个递减序列 \((\mathbf{A}_n)_{n\geqslant 1}\)（由 \(\Phi\) 的元素组成、满足 \(\bigcap_{n\geqslant 1}\mathbf{A}_n = \emptyset\)）成立。若 \(m\) 可数可加，证明 \(|m|\) 也可数可加。

\begin{enumerate}
\def\labelenumi{\arabic{enumi})}
\setcounter{enumi}{1}
\tightlist
\item
  设 E 为一个集合，\(\Phi\) 为 E 的子集所成的一个集环，V 为一个全格序的 Riesz 空间。称映射 \(m\)（\(\Phi\) 到 V 中的）为正的，如果有 \(m(A) \geqslant 0\)（对所有 \(A \in \Phi\)）；称其为可加的，如果有 \(m(A \cup B) = m(A) + m(B)\)（当集合 \(A \in \Phi\) 与 \(B \in \Phi\) 不相交时）；称其为相对有界的，如果元素 \(m(B)\)（其中 \(B \in \Phi\)，\(B \subset A\)）所成的集合在 V 中有上界（对任何 \(A \in \Phi\)）。
\end{enumerate}

\begin{enumerate}
\def\labelenumi{\alph{enumi})}
\tightlist
\item
  设 \(m\) 是 \(\Phi\) 到 \(V\) 中的一个相对有界映射。对每个 \(A\)，按 Exer. 1 a) 的方式定义元素 \(m_{+}(A), m_{-}(A)\) 和 \(|m|(A)\)（均属于 \(V\)）。证明 \(|m|(A)\)
\end{enumerate}

是 V 中形如 \(\sum_{i=1}^{p}|m(A_i)|\) 的元素所成集合的上确界（对所有有限划分 \((\mathbf{A}_i)_{1\leqslant i\leqslant p}\)（把 A 划分成属于 \(\Phi\) 的集合））。由此推出映射 \(|m|\)（\(\Phi\) 到 V 中的）是可加的，继而推出 \(m\) 是两个可加正映射 \(m_{+}\) 与 \(m_{-}\)（均由 \(\Phi\) 到 V 中）之差。

\begin{enumerate}
\def\labelenumi{\alph{enumi})}
\setcounter{enumi}{1}
\tightlist
\item
  设 \(m'\) 与 \(m''\) 是 \(\Phi\) 到 \(V\) 中的两个可加正映射，且 \(m = m' - m''\)。证明 \(m\) 相对有界，且 \(m' \geqslant m_+\)，\(m'' \geqslant m_-\)。c) 推广 Exer. 1 d)。
\end{enumerate}

\begin{enumerate}
\def\labelenumi{\arabic{enumi})}
\setcounter{enumi}{2}
\tightlist
\item
  设 \(\mathbf{E}\) 为一个集合，\(\mathfrak{T}\) 为 \(\mathbf{E}\) 的子集所成的一个 σ-集环。用 \(\mathcal{M}\) 表示映射 \(f\)（\(\mathbf{E}\) 到 \(\overline{\mathbf{R}}_+\) 中的）所成的集合，这些映射使得 \(x \in \mathbf{E}\) 中使 \(f(x) \geqslant c\) 的 x 所成的集合属于 \(\mathfrak{T}\)（对每个 \(c \in \overline{\mathbf{R}}_+\)）。
\end{enumerate}

\begin{enumerate}
\def\labelenumi{\alph{enumi})}
\item
  证明 \(\mathcal{M}\) 中元素组成的每个序列的上极限和下极限都属于 \(\mathcal{M}\)（先研究递增或递减的序列）。
\item
  证明 \(\mathfrak{T}\) 是 \(\mathbf{E}\) 中这样的子集所成的集合：其特征函数属于 \(\mathcal{M}\)。
\item
  设 \(f \in \mathcal{M}\)。对每个 \(B\)（\(\overline{\mathbf{R}}_+\) 的 Borel 子集），集合 \(f^{-1}(B)\) 属于 \(\mathfrak{T}\)（子集 \(B\)（属于 \(\overline{\mathbf{R}}_+\)）中使 \(f^{-1}(B) \in \mathfrak{T}\) 的那些 B 所成的集合是一个 σ-集环，且包含区间 \([c, +\infty]\)）。
\item
  设 \(f_1, \ldots, f_n\) 属于 \(\mathcal{M}\)，\(\varphi\) 是 \((\overline{\mathbf{R}}_+)^n\) 到 \(\overline{\mathbf{R}}_+\) 中的一个 Borel 映射；证明映射 \(x \mapsto \varphi(f_1(x), \ldots, f_n(x))\)（E 到 \(\overline{\mathbf{R}}_+\) 中的）属于 \(\mathcal{M}\)。由此推出 \(\mathcal{M}\) 包含由 \(\mathcal{M}\) 的元素组成的每个级数的和。
\item
  证明 \(\mathcal{M}\) 是有限正 \(\mathfrak{T}\) -阶梯函数组成的递增序列的极限所成的集合。
\end{enumerate}

\begin{enumerate}
\def\labelenumi{\arabic{enumi})}
\setcounter{enumi}{3}
\tightlist
\item
  记号与假设同上一习题。所谓 \((\mathbf{E},\mathfrak{T})\) 上的抽象测度，是指每个映射 \(m\)（从 \(\mathfrak{T}\) 到 \(\overline{\mathbf{R}}_+\) 中，且具有下列性质）：对每个序列 \((\mathrm{A}_n)_{n\in \mathbf{N}}\)（由属于 \(\mathfrak{T}\) 的集合组成，且两两不相交），有 \(m\big(\bigcup_{n\in \mathbf{N}}\mathrm{A}_n\big) =\)
\end{enumerate}

\(\sum_{n\in \mathbf{N}}m(\mathrm{A}_n)\)。所谓 \((\operatorname {E},\mathfrak{T})\) 上的积分，是指 \(\mathcal{M}\) 到 \(\overline{\mathbf{R}}_+\) 中且满足下列条件的每个映射 J：

  下列条件：\(\alpha)\) \(J(\lambda \cdot f) = \lambda \cdot J(f)\) 对 \(\lambda \in \overline{\mathbf{R}}_{+}\) 与 \(f \in \mathcal{M}\) 成立（按通常的约定 \(0 \cdot (+\infty) = 0\)）；\(\beta)\) \(J\big(\sum_{n \in \mathbf{N}} f_n\big) = \sum_{n \in \mathbf{N}} J(f_n)\) 对每个序列 \((f_n)_{n \in \mathbf{N}}\)（由 \(\mathcal{M}\) 的元素组成）成立。


\begin{enumerate}
\def\labelenumi{\alph{enumi})}
\item
  设 \(J\) 是 \((\mathbf{E},\mathfrak{T})\) 上的一个积分；对每个子集 \(A \in \mathfrak{T}\)，置 \(m_{J}(A) = J(\varphi_A)\)；证明 \(m_J\) 是一个抽象测度，且映射 \(J \mapsto m_J\) 是 \((\mathbf{E},\mathfrak{T})\) 上积分所成的集合到 \((\mathbf{E},\mathfrak{T})\) 上抽象测度所成的集合上的一个双射（若 \(m\) 是抽象测度，先定义 \(J(f)\)（按线性，对有限正 \(\mathfrak{T}\) -阶梯函数），并利用 Exer. 3 e)）。若 \(m\) 是抽象测度，我们将用 \(m^*\) 表示对应的积分。
\item
  设 \(m\) 是 (E, T) 上的一个抽象测度。对 E 上的每个正函数 \(f\)（无论有限与否），置 \(m^{*}(f) = \inf_{g \in \mathcal{M}, g \geqslant f} m^{*}(g)\)。证明 \(m^{*}\) 是 E 上的一个负担（模仿 Ch. IV, §1, No.~3 的 Th. 3 的证明）。对 E 的每个子集 A，置 \(m^{*}(A) = m^{*}(\varphi_{A})\)。
\end{enumerate}

\begin{enumerate}
\def\labelenumi{\arabic{enumi})}
\setcounter{enumi}{4}
\tightlist
\item
  设 \(\mathbf{E}\) 为一个集合，\(\mathfrak{T}\) 为 \(\mathbf{E}\) 的子集所成的一个 σ-集环，\(m\) 为 \((\mathbf{E},\mathfrak{T})\) 上的一个抽象测度，\(p\geqslant 1\) 为一个有限实数。对每个数值函数 \(f\)（\(\mathbf{E}\) 上的），置 \(\mathrm{N}_p(f) = m^* (|f|^p)^{1 / p}\)，并用 \(\mathcal{F}^p\) 表示函数 \(f\)（使 \(\mathrm{N}_p(f)\) 为有限）所成的集合（cf.~§1, Exer. 4）。给向量空间 \(\mathcal{F}^p\) 赋以半范数 \(\mathrm{N}_p\)，对它而言该空间是完备的。
\end{enumerate}

\begin{enumerate}
\def\labelenumi{\alph{enumi})}
\item
  设 \(\mathcal{N}\) 为使 \(m^{*}(|f|)=0\) 的函数 f 所成的集合，\(\mathcal{V}\) 为由属于 \(\mathcal{M}\) 的有限函数生成的向量空间，\(\mathcal{L}^{p}\) 为 \(\mathcal{F}^{p}\) 中包含集合 \(A\in \mathfrak{T}\)（使 \(m^{*}(A)\) 为有限）的特征函数的最小闭线性子空间。证明 \(\mathcal{L}^{p}=\mathcal{N}+(\mathcal{V}\cap\mathcal{F}^{p})\)。
\item
  把 Lebesgue 定理（Ch. IV, §3, No.~7, Th. 6）推广到目前的情形，并特别考察 \(p = 1\) 的情形（此时应定义函数 \(f \in \mathcal{L}^1\) 的积分）。
\item
  称 \(\mathbf{A}\)（\(\mathbf{E}\) 的一个子集）是 \(m\) -可积的，如果 \(\varphi_{\mathbf{A}} \in \mathcal{L}^1\)。证明此时存在集合 \(\mathbf{A}' \in \mathfrak{T}\)，使得 \(m^{*}(\mathbf{A} \cap \mathbb{C}\mathbf{A}') = m^{*}(\mathbf{A}' \cap \mathbb{C}\mathbf{A}) = 0\)。在 \(m(\mathbf{E})\) 有限的情形，逆命题成立。
\end{enumerate}

¶ 6) 设 E 为一个集合，T 为 E 的子集所成的一个 σ-集环，m 为 (E, T) 上的一个抽象测度。假设 \(m(E)\) 有限，并且对每个 \(B \in T\)（满足 \(m(B) > 0\)），存在集合 \(A \in T\)，使得 \(A \subset B\) 且 \(0 < m(A) < m(B)\)。则对每个 \(B \in T\) 与每个满足 \(0 \leqslant t \leqslant m(B)\) 的数 t，存在集合 \(A \in T\)，使得 \(m(A) = t\) 且 \(A \subset B\)。

¶ 7) 设 T 为一个拓扑空间，\(\Phi\) 为 T 的子集所成的一个集环，\(m\) 为 \(\Phi\) 到 R 中的一个相对有界可加映射（cf.~Exer. 1）。假设存在 T 的开覆盖 \(\mathfrak{U}\)，使得 \(\Phi\) 由包含在 \(\mathfrak{U}\) 中有限多个元素的并内的 T 的 Borel 子集组成。再假设对每个 \(A \in \Phi\) 与每个 \(\varepsilon > 0\)，存在 T 的紧子集 K 与 T 的开子集 U，使得 \(K \subset A \subset U\)，\(U \in \Phi\) 且 \(|m|(U - K) < \varepsilon\)。证明 T 上存在测度 \(\mu\)（不必为正），使得 \(m(A) = \mu^{\bullet}(A)\)（对所有 \(A \in \Phi\)），并且这样的测度 \(\mu\) 是唯一的（通过引入 Exer. 1 c) 的分解 \(m = m_{+} - m_{-}\)，化为 \(m\) 取正值的情形；先处理 \(T \in \Phi\) 的情形，并注意此时 \(m\) 是内正则的；最后通过把测度粘合起来处理一般情形）。

\begin{enumerate}
\def\labelenumi{\arabic{enumi})}
\setcounter{enumi}{7}
\tightlist
\item
  设 T 为一个拓扑空间，m 为 \(\mathfrak{B}(\mathrm{T})\) 到 R 中的一个有界、可数可加的映射（cf.~Exer. 1d)）。用 T 表示具有下列性质的 T 的 Borel 子集 A 所成的集合：对每个 \(\varepsilon > 0\)，存在 T 中的闭集 F 与开集 U，使得 \(F \subset A \subset U\)，并且对 U - F 的每个 Borel 子集 B 有 \(|m(B)| < \varepsilon\)（换句话说，\(|m|(\mathrm{U} - \mathrm{F}) < \varepsilon\)）。证明 T 是 T 的子集所成的一个 σ-集环（先注意 \(T \in T\)，且 T 是一个集环；然后只须证明：若 \((\mathrm{A}_{n})_{n \in \mathbb{N}}\) 是 T 中两两不相交的元素组成的序列，则集合 \(A = \bigcup_{n \in N} A_{n}\) 属于 T；为此，选取闭集 \(F_{n}\) 与开集 \(U_{n}\)，使得 \(F_{n} \subset A_{n} \subset U_{n}\) 且 \(|m|(\mathrm{U}_{n} - \mathrm{F}_{n}) < \varepsilon / 2^{n}\)，然后对充分大的 p 置 \(F = F_{0} \cup \cdots \cup F_{p}\)，并置 \(U = \bigcup_{n \in \mathbf{N}} U_{n}\)。
\end{enumerate}


\begin{enumerate}
\def\labelenumi{\arabic{enumi})}
\setcounter{enumi}{8}
\tightlist
\item
  设 \(\mathbf{X}\) 为一个集合；称 \(\mathbf{X}\) 上的每个可数可加映射（\(\mathfrak{P}(\mathbf{X})\) 到 \(\mathbf{R}_+\) 中的）为负担。若负担 \(m\)（\(\mathbf{X}\) 上的）满足 \(m(\{x\}) = 0\)（对所有 \(x \in \mathbf{X}\)），则称它是弥散的；若存在正函数 \(f\)（在 \(\mathbf{X}\) 上），使得 \(m(\mathbf{A}) = \sum_{x \in \mathbf{A}} f(x)\) 对每个
\end{enumerate}

X 的子集 A 都成立，则称其为原子的。

\begin{enumerate}
\def\labelenumi{\alph{enumi})}
\item
  证明 X 上的每个负担都可以以唯一的方式分解为一个弥散负担与一个原子负担之和。
\item
  设 \(X'\) 为 \(X\) 的一个子集，\(m'\) 为 \(X'\) 上的一个负担；对每个子集 \(A\)（\(X\) 的），置 \(m(A) = m'(A \cap X')\)。证明 \(m\) 是 \(X\) 上的一个负担，并且它是弥散的（相应地，原子的）当且仅当 \(m'\) 具有同样的性质。
\item
  设 \(m\) 为 \(X\) 上的一个负担，\((X_i)_{i \in I}\) 为 \(X\) 的两两不相交的子集组成的一个族。若 \(I\) 上的每个负担都是原子的，则 \(m\big(\bigcup_{i \in I} X_i\big) = \sum_{i \in I} m(X_i)\)。
\end{enumerate}

\begin{enumerate}
\def\labelenumi{\arabic{enumi})}
\setcounter{enumi}{9}
\tightlist
\item
  设 \(\mathbf{X}\) 为一个不可数的无限集，带有一个良序关系，记作 \(x \leqslant y\)。对每个 \(x \in \mathbf{X}\)，用 \(\mathrm{I}(x)\) 表示 \(y \in \mathbf{X}\) 中满足 \(y < x\) 的 y 所成的集合。我们作下列假设：\(\alpha\) ) 存在一个最大元素 \(a\)（在 \(\mathbf{X}\) 中）；\(\beta\) ) 严格小于 \(\operatorname{Card}(\mathbf{X})\) 的基数所成的集合有一个最大元素 \(\mathfrak{c}\)；\(\gamma\) ) 对每个 \(x\)（在 \(\mathrm{I}(a)\) 中），有 \(\operatorname{Card}(\mathrm{I}(x)) \leqslant \mathfrak{c}\)，且集合 \(\mathbf{Y}\)（由 \(x \in \mathbf{X}\) 中使 \(\operatorname{Card}(\mathrm{I}(x)) = \mathfrak{c}\) 的 x 组成）非空；
\end{enumerate}

\(\delta)\) 基数 \(\leqslant \mathfrak{c}\) 的集合上的每个负担都是原子的。证明 X 上的每个负担都是原子的。可以按如下方式用反证法论证：设 \(m\) 为 X 上的一个非零弥散负担。用 \(b\) 表示 Y 的最小元素，并对每个 \(x \in I(a)\)，令 \(f_x\) 为 \(I(x)\) 到 \(I(b)\) 中的一个单射。对每对 \((x,y) \in I(a) \times I(b)\)，用 \(A_{x,y}\) 表示由元素 \(z \in X\)（满足 \(x < z < a\) 且 \(f_z(x) = y\)）所成的集合；对每个 \(x \in I(a)\)，用 \(M_x\) 表示由 \(y \in I(b)\)（满足 \(m(A_{x,y}) > 0\)）所成的集合，并对每个 \(y \in I(b)\)，用 \(N_y\) 表示由 \(x \in I(a)\)（满足 \(m(A_{x,y}) > 0\)）所成的集合。利用 Exer. 9 c)，证明 \(m(X) = \sum_{y \in I(b)} m(A_{x,y})\)；由此推出集合 \(M_x\) 对每个 \(x \in I(a)\) 都是可数的且非空；再证明每个集合 \(N_y\) 都是可数的，并由此得出矛盾 \(\text{Card}(X) = c\)。

\begin{enumerate}
\def\labelenumi{\arabic{enumi})}
\setcounter{enumi}{10}
\tightlist
\item
  称基数 \(\mathfrak{c}\) 为 Ulam 基数（或 Ulam 型的），如果每个具有基数 \(\mathfrak{c}\) 的集合上的负担都是原子的。证明每个可数基数都是 Ulam 型的；若 \(\mathfrak{c}\) 是 Ulam 基数，则每个基数 \(\mathfrak{c}' \leqslant \mathfrak{c}\) 也是 Ulam 基数；并且若 \((\mathfrak{c}_i)_{i \in I}\) 是一族 Ulam 基数，且 \(\operatorname{Card}(I)\) 是 Ulam 型的，则基数 \(\sum_{i \in I} \mathfrak{c}_i\) 是 Ulam 型的。最后，
\end{enumerate}

最小不可数基数是 Ulam 型的（cf.~Exer. 10）。

\begin{enumerate}
\def\labelenumi{\arabic{enumi})}
\setcounter{enumi}{11}
\tightlist
\item
  设 X 为一个集合，U 为 \(\mathfrak{P}(X)\) 的一个子集，m 为 U 的特征函数。为使 m 是一个负担，必要且充分的是 U 是一个超滤子，并且 U 的每个可数子集的交属于 U。在这些条件下，称 U 为 X 上的一个 Ulam 超滤子。假设情形如此，且 \((\mathrm{H}_{i})_{i\in\mathrm{I}}\) 是 U 中元素组成的一个族，其中 \(\operatorname{Card}(\mathrm{I})\) 是 Ulam 基数；证明 \(\bigcap_{i\in I} H_{i}\) 属于
\end{enumerate}

\(\mathfrak{U}\)。

\begin{enumerate}
\def\labelenumi{\arabic{enumi})}
\setcounter{enumi}{12}
\tightlist
\item
  在本习题中，我们承认连续统假设，也就是说，我们把公理 \(\langle\mathrm{Card}(\mathbf{R})\) 是最小不可数基数 \(\rangle\) 附加到集合论的公理中。\(^{(1)}\)
\end{enumerate}

\begin{enumerate}
\def\labelenumi{\alph{enumi})}
\item
  \(\mathbf{R}\) 上的每个负担都是原子的（cf.~Exer. 11）。
\item
  设 X 为一个集合，m 为 X 上满足下列性质的负担：对每个满足 \(m(A) > 0\) 的 X 的子集 A，存在 A 的子集 B，使得 \(0 < m(B) < m(A)\)；则 m 为零（对每个整数 \(n \geqslant 1\)，存在 X 的一个有限划分 \((\mathrm{X}_{n,i})_{i \in \mathrm{I}_n}\)，使得 \(m(\mathrm{X}_{n,i}) \leqslant 1/n\)（对所有 \(i \in I_n\)）；置 \(I = \prod_{n \geqslant 1} I_n\) 与
\end{enumerate}

\(\mathrm{X}_{\alpha} = \bigcap_{n\geqslant 1}\mathrm{X}_{n,\alpha_n}\)（对 I 中的每个 \(\alpha = (\alpha_{n})_{n\geqslant 1}\)）；族 \((\mathrm{X}_i)_{i\in \mathrm{I}}\) 是 \(\mathbf{X}\) 的一个划分，

\(m(X_{i}) = 0\)（对每个 \(i \in I\)），并且 \(I\) 上的每个负担由 \(a\) ) 知都是原子的；借助 Exer. 9 c) 完成证明。

\begin{enumerate}
\def\labelenumi{\alph{enumi})}
\setcounter{enumi}{2}
\item
  设 X 是一个不存在非平凡 Ulam 超滤子的集合；则 X 上的每个负担都是原子的（否则，由 b)，X 上将存在一个弥散负担 m 与 X 的子集 A，使得 \(m(A) > 0\)，且对 A 的每个子集 B，\(m(B) = 0\) 或 \(m(A - B) = 0\)；由满足 \(m(A) = m(A \cap B)\) 的 X 的子集 B 组成的集合 U 将是 X 上的一个非平凡 Ulam 超滤子）。
\item
  证明 \(2^{\mathfrak{c}}\) 对每个 Ulam 基数 \(\mathfrak{c}\) 都是 Ulam 基数（由 \(c\) )，只须证明：若 \(C\) 是具有基数 \(\mathfrak{c}\) 的集合，则每个 Ulam 超滤子 \(\mathfrak{U}\)（在 \(X = \{0,1\}^{C}\) 上）都是平凡的；给 \(C\) 装备一个良序结构，并用超限归纳定义一个族 \(s = (s_{\alpha})_{\alpha \in C}\)（由 \(\{0,1\}\) 中元素组成），使得由 \(x = (x_{\alpha})_{\alpha \in C}\)（\(X\) 中满足 \(x_{\alpha}=s_{\alpha}\)（对所有 \(\alpha\leqslant\beta\)）者）组成的集合属于 \(\mathfrak{U}\)，对每个 \(\beta\in\mathrm{C}\)（cf.~Exer. 12）；于是 \(s\) 属于 \(\mathfrak{U}\) 的每个元素）。\(^{(2)}\)
\end{enumerate}

\begin{enumerate}
\def\labelenumi{\arabic{enumi})}
\setcounter{enumi}{13}
\item
  设 \((\mathbf{T}_i)_{i\in I}\) 是一族 Radon 拓扑空间，\(\mathbf{T}\) 是该族的和空间。若 \(\operatorname{Card}(\mathbf{I})\) 是 Ulam 基数，则拓扑空间 \(\mathbf{T}\) 是 Radon 空间（证明：对每个正的、有界的可数可加函数 \(m\)（在 \(\mathbf{T}\) 的 Borel σ-集环上），存在可数子集 \(J\)（\(I\) 的），使得 \(m(\bigcup_{i \in I - J} \mathbf{T}_i) = 0\)）。
\item
  设 T 是一个离散空间，其基数是最小不可数基数。证明 T 是一个 Radon 局部紧空间，其拓扑没有可数基（cf.~Exer. 11）。由此推出存在不可度量化的紧 Radon 空间。
\item
  设 \(I\) 是一个可数集，\((\mathbf{T}_i)_{i \in I}\) 是一族 Radon 空间。假设这些空间 \(\mathbf{T}_i\) 中每一个的紧子集都是可度量化的。证明乘积空间 \(\mathbf{T} = \prod_{i \in I} \mathbf{T}_i\) 是 Radon 空间。推广到逆极限的情形。
\item
  设 \(\mathbf{T}\) 是一个拓扑空间。
\end{enumerate}

\begin{enumerate}
\def\labelenumi{\alph{enumi})}
\tightlist
\item
  设 \(m\) 是 \(\mathfrak{B}(\mathbf{T})\) 到 \(\mathbf{R}_+\) 中的一个可数可加映射。假设存在一个 Borel 子集的序列 \((\mathbf{T}_n)_{n\in \mathbf{N}}\)（\(\mathbf{T}\) 的），使得 \(\mathbf{T} = \bigcup_{n\in \mathbf{N}}\mathbf{T}_n\)，\(m(\mathbf{T}_0) = 0\)
\end{enumerate}

并且子空间 \(\mathbf{T}_n\)（\(\mathbf{T}\) 的）对每个 \(n \geqslant 1\) 都是 Radon 的。证明存在一个有界正测度 \(\mu\)（在 \(\mathbf{T}\) 上），使得 \(m(\mathbf{A}) = \mu^{\bullet}(\mathbf{A})\)（对每个 Borel 子集 \(\mathbf{A}\)（\(\mathbf{T}\) 的））（可以借助 No.~3 的 Prop. 2 的推论，化为这些 \(\mathbf{T}_n\) 两两不相交的情形）。

\begin{enumerate}
\def\labelenumi{\alph{enumi})}
\setcounter{enumi}{1}
\item
  把上一结果推广到如下情形：集合 \(T_{n}\)（对 \(n \geqslant 1\)）不再假设是 Borel 的，而只是 T 中普遍可测的（按 No.~3 的 Prop. 2 的方式论证）。
\item
  每个是一个可度量化紧子空间序列之并的拓扑空间都是 Radon 的。
\end{enumerate}

\begin{enumerate}
\def\labelenumi{\arabic{enumi})}
\setcounter{enumi}{17}
\tightlist
\item
  设 \(\mathbf{T}_1\) 与 \(\mathbf{T}_2\) 是两个拓扑空间，\(f\) 是 \(\mathbf{T}_1\) 到 \(\mathbf{T}_2\) 上的一个连续双射映射。假设对每个有界可数可加函数 \(m_1\)（在 Borel σ-集环 \(\mathfrak{B}(\mathbf{T}_1)\) 上），存在一个由紧子集 \(\mathrm{K}_n\)（\(\mathbf{T}_1\) 的）组成的序列，使得 \(m_1(\mathbf{T}_1 - \bigcup_{n}\mathbf{K}_n) = 0\)。设 \(m_2\) 是 \(\mathfrak{B}(\mathbf{T}_2)\) 上的一个有界可数可加函数。
\end{enumerate}

为使存在一个有界可数可加函数 \(m_{1}\)（在 \(\mathfrak{B}(\mathrm{T}_{1})\) 上），使得 \(m_{2}(\mathrm{A}) = m_{1}(f^{-1}(\mathrm{A}))\)（对所有 \(\mathrm{A} \in \mathfrak{B}(\mathrm{T}_{2})\)），必要且充分的是下列条件成立：存在一个由紧子集 \(K_{n}\)（\(T_{1}\) 的）组成的序列，使得 \(m_{2}(\mathrm{T}_{2} - \bigcup_{n} f(\mathrm{K}_{n})) = 0\)。

